\section{Deliberately deferred / not adopted}\label{sec:deferred}

A design is defined as much by what it refuses as by what it adopts. The station-pair
backbone of \cref{sec:datamodel} was selected from five competing sketches and hardened
against an adversarial critique; that process produced a set of alternatives that were
explicitly considered and set aside. This section records each one together with the
reasoning, so that a future contributor revisiting any of these choices inherits the
argument rather than rediscovering it. The rejections fall into two kinds: choices that
are \emph{wrong} for QtRocket and should stay rejected, and choices that are \emph{correct
in the limit} but premature for a rigid inline stack and are therefore deferred rather than
dismissed. The distinction matters --- a deferred idea is a planned extension point, not a
dead end.

\begin{description}

  \item[\zaft\ as the longitudinal convention --- \textbf{rejected (sign-inverted).}]
  Adopting \zaft\ would invert the longitudinal axis relative to every component that
  already exists. The visualizer builds its primitives with the fore extent at positive
  $z$ (\code{buildCone} and \code{buildTube} place the tip / top at $+z$), and every legacy
  \code{.qrd} fixture --- including \src{tests/data/designs/xl75\_multi.qrd}{1} --- expresses
  its offsets in a \zfwd\ frame. Choosing \zaft\ would force a whole-axis inversion of the
  renderer (cap normals, sphere poles, every primitive) \emph{and} a sign-negating
  migration of every stored file. Adopting \zfwd\ instead reduces the renderer change to a
  consumer swap with no sign inversion, and the file migration to a pure re-expression with
  no axial flip. The convention was settled once, in agreement with the code and the data
  on disk; see \cref{sec:philosophy}.

  \item[A stored resolved transform as ground truth --- \textbf{rejected (reintroduces
  staleness).}] One tempting simplification is to resolve each part's absolute pose once
  and store it as the canonical placement. This was rejected because it reintroduces a
  staleness class that is the exact inverse of the bug this design exists to remove: edit a
  part's length and its stored transform is silently wrong, with nothing to recompute it.
  QtRocket already commits to a ``derive, don't store'' discipline for composite
  quantities, and this design extends that discipline to placement. The \code{Pose} /
  \code{compose} mechanics and the snap-operator verbs are adopted from the scene-graph
  sketch, but only \emph{intent} is stored; the absolute pose is re-derived on every read.

  \item[A rich named-datum enum with an escape hatch --- \textbf{rejected
  (over-engineering).}] A datum vocabulary such as \cppinline|enum { FORE, MID, AFT, ... }|
  with a parametric escape hatch for arbitrary positions was rejected as paying two costs at
  once: the rigidity of a closed enum \emph{and} the complexity of the mechanism for
  escaping it. The named edges are merely fractions --- \code{FORE}, \code{MID}, and
  \code{AFT} are exactly $1$, $0.5$, and $0$ --- and an arbitrary station such as a
  mid-tube rail-button location is just another fraction in $[0,1]$. Two fractional
  stations plus a small closed \code{SeatKind} therefore express everything the enum would,
  with nothing to extend (\cref{sec:datamodel}).

  \item[Parallel \code{childParts} and \code{childLinks} vectors --- \textbf{rejected
  (desync hazard).}] Storing the child pointers and their links in two side-by-side vectors
  was rejected because the two arrays must be kept index-aligned by hand across every
  structural mutation. The single-vector storage at \src{model/parts/Part.h}{326} is
  touched by \code{clone()}, \code{removeChildById}, and \code{addChildPart}; splitting it
  invites an off-by-one or a missed erase that silently attaches the wrong link to the
  wrong child. The design keeps a single vector and changes only the element type, from a
  \cppinline|std::tuple<std::shared_ptr<Part>, Vector3>| to a
  \cppinline|std::pair<std::shared_ptr<Part>, StationLink>| (\cref{sec:datamodel}), so the
  pointer and its link cannot drift apart.

  \item[``Every link mandatory'' --- \textbf{rejected (punishes the trivial stack).}]
  Requiring an explicit link on every attachment was rejected because it taxes the common
  case to no benefit. The overwhelmingly typical relationship --- a child abutting aft of
  its parent in a straight inline stack --- should cost zero authored numbers. The
  defaulted \seat{Abut} link (parent aft plane to child fore plane, zero gap) preserves
  that, so \cppinline|nose->addChildPart(body)| remains a valid, fully specified
  attachment; see \cref{sec:authoring}.

  \item[Scalar \code{radiusOuterAt} for \code{FinSet} --- \textbf{rejected (dishonest).}]
  A fin set is not axisymmetric, so reporting a single scalar outer radius for the overlap
  sweep would be a fabrication: a \cppinline|bodyRadius + span| disc would collide falsely
  with every adjacent body tube. The design instead has a fin set contribute only its
  \emph{body disc} to the axial envelope, with true azimuthal (sectored) fin interference
  scoped out of v1 and documented as unchecked, pending a fin-clocking model QtRocket does
  not yet have. This is honest under-reporting rather than confident wrong-reporting; the
  diagnostics treatment is detailed in \cref{sec:diagnostics}.

  \item[Off-axis ($r > 0$) placement in 3-DOF --- \textbf{deferred to 6-DOF.}] Placing a
  part radially off the longitudinal axis --- a single rail button, an asymmetric pod --- is
  deferred because seating mass at $x = r$ without an accompanying orientation is itself a
  half-correct stored placement, the very error class this design eliminates. The radial
  value $r$ is therefore used \emph{now} only for the seam and overlap \emph{check}, while
  the pose solve sets $x = y = 0$; correct off-axis seating arrives together with
  \code{childRot} and a per-part body-frame CM offset in the 6-DOF work (\cref{sec:sixdof}).

\end{description}

Three of these warrant fuller treatment, because each guards a correctness invariant that
is easy to violate under future pressure to ``simplify.''

\begin{rejected}[Storing the resolved transform]
The single most attractive --- and most dangerous --- shortcut is to cache the absolute
\code{Pose} of each part as its source of truth. It would make reads trivially cheap and
the data model flat. It is rejected because it merely relocates the staleness bug from the
input side to the output side: today a stale \emph{CM-to-CM offset} disagrees with geometry
when a length is edited; a stored \emph{resolved transform} would disagree with geometry in
exactly the same situations, with no mechanism short of a full re-resolve to detect it. The
design keeps placement strictly derived and gates the re-derive on a structural-dirty flag
(\cref{sec:resolver}), so the cache is an optimization that can never be wrong, not a second
copy of the truth that can.
\end{rejected}

\begin{rejected}[A rich named-datum enum with an escape hatch]
A vocabulary of named datums looks more expressive than two bare fractions, but it is
strictly worse. It hard-codes the three edges (\code{FORE}/\code{MID}/\code{AFT}) that are
already representable as $1$/$0.5$/$0$, and then needs a parametric escape hatch for every
station that is not one of those edges --- so it pays for a closed enum and for the machinery
to open it. The chosen representation, two stations in $[0,1]$ plus a closed three-value
\code{SeatKind}, is both smaller and more general: there is no enum to extend because there
is nothing the enum was buying. \code{SeatKind} itself stays closed deliberately, because a
rigid axisymmetric stack genuinely expresses only three radial relationships (\cref{sec:datamodel}).
\end{rejected}

The most consequential ``not yet'' is the full joint graph. Unlike the rejections above, it
is the correct long-term model and is deferred only because the present scope does not need
it.

\begin{rejected}[A full joint graph with a sibling \texttt{<joints>} block]
A general design would model attachments as a graph of joint entities, serialized in a
sibling \mintinline{xml}{<joints>} block keyed independently of the part tree. This is the
correct long-term chassis: clusters, struts, off-axis pods, and the 6-DOF
companion-concentricity constraint that \cref{sec:diagnostics}'s Layer 3 will eventually need
all require a part to relate to more than its single tree parent --- relationships an
ownership tree cannot express. It is nonetheless deferred, for two reasons. First, for a
rigid inline stack it is over-engineering: every relationship QtRocket authors today is a
single parent-to-child seam, which the inline \cppinline|std::pair<std::shared_ptr<Part>, StationLink>|
captures exactly. Second, a sibling joints block keyed by name is an active hazard given
that part names need not be unique (\src{model/parts/Part.h}{193}); keeping the link
\emph{inline} with its part sidesteps that ambiguity entirely (\cref{sec:serialization}).
The inline pair is forward-compatible with the graph: when real multi-host designs appear,
the single \code{StationLink} per child can widen into a list of links --- or a separate
joint table --- without disturbing the resolver's interface or the serialization of the
common case.
\end{rejected}

\begin{keyinsight}[Deferred is not dismissed]
The joint graph and off-axis seating are not rejected ideas; they are staged ones. The data
model is deliberately shaped so that adopting them later is additive: \code{Pose} already
carries a \code{Quaternion} and \code{compose()} already rotates child translations
correctly (the ``no rotation'' limit at \src{model/parts/Part.h}{43} is a property of
today's identity \code{childRot}, not of the resolver), and the single-link-per-child
storage can widen to a graph. Each deferral therefore leaves a defined extension point
rather than a wall, which is why the convention, the seat semantics, and the \code{Pose}
mechanics were all settled now even though only their 3-DOF subset is exercised today.
\end{keyinsight}
